\chapter{अंडाणु से जन्म तक}\label{ch:g8:from-egg-to-birth}

\Cref{ch:g5:human-reproduction} ने उन नौ महीनों को एक कहानी की तरह
सुनाया था; और अब तुम्हारे पास वे सारे साज़ हैं जो उस कहानी को ठीक से
बजाने के लिए चाहिए: \omterm{def:g5:first-look-at-cells:cell}{कोशिकाएँ} और उनका बँटना, दोनों \omterm{def:g8:sexual-reproduction:gametes}{युग्मक} और उनकी
मिलन-जगह, \omterm{def:g8:puberty:hormones}{हॉर्मोन}, और \omterm{prop:g4:journey-of-food:absorb}{रक्त} की पहुँचाई हुई चीज़ें। यह अध्याय इंसानी विकास
को मशीनरी खोलकर दोबारा बजाता है --- एक \omterm{def:g5:first-look-at-cells:cell}{कोशिका} से तीन हज़ार अरब तक,
चालीस हफ़्तों में।

\section{पहला हफ़्ता: युग्मनज से किरायेदार तक}

\begin{proposition}[अंडवाहिनी में नीचे का सफ़र]\label{prop:g8:from-egg-to-birth:firstweek}
\omterm{def:g8:reproductive-systems:female}{अंडवाहिनी} में \omterm{def:g5:flower-to-fruit:steps}{निषेचन}
(\cref{def:g8:reproductive-systems:female}) एक ही साथ दो घड़ियाँ चालू
कर देता है:
\begin{enumerate}
  \item \emph{बँटना}: \omterm{prop:g8:sexual-reproduction:core}{युग्मनज} एक दिन के भीतर बँटता है, फिर बार-बार
        (\cref{ex:g5:first-look-at-cells:one}) --- दो, चार, आठ
        \omterm{def:g5:first-look-at-cells:cell}{कोशिकाएँ}, और चौथे दिन तक दर्जनों की एक शहतूत जैसी गेंद, और यह
        सब \omterm{def:g8:reproductive-systems:female}{अंडवाहिनी} में \omterm{prop:g5:human-reproduction:plan}{गर्भाशय} की ओर बहते-बहते;
  \item \emph{आरोपण}: सातवें दिन के आस-पास वह नन्ही गेंद \omterm{prop:g5:human-reproduction:plan}{गर्भाशय} तक
        पहुँचती है और तैयार परत में ख़ुद को धँसा लेती है
        (\cref{prop:g8:reproductive-systems:cycle} की चौथी घटना) ---
        किरायेदार आ पहुँचा, चक्र रुक गया, और \omterm{def:g8:puberty:hormones}{हॉर्मोन} पूरे तंत्र को
        इशारा कर देते हैं: \omterm{def:g5:human-reproduction:pregnancy}{गर्भावस्था}।
\end{enumerate}
\end{proposition}
\admitted

\begin{example}[पहले बँटवारे जीते किस पर हैं]\label{ex:g8:from-egg-to-birth:provisions}
पहले पूरे हफ़्ते वह मुसाफ़िर न कोई खाना खाता है न उसे कुछ पहुँचाया जाता
है: वह ख़ुद \omterm{def:g4:animal-reproduction:sexual}{अंडाणु} के साज़-सामान पर चलता है --- यानी ठीक उसी चीज़ पर,
जिसके लिए \cref{ex:g8:reproductive-systems:egg} वाली भरपाई थी। \omterm{def:g8:reproductive-systems:male}{शुक्राणु}
पिता का \omterm{def:g6:cell-unit-of-life:parts}{केंद्रक} लाया था और उसके सिवा कुछ ख़ास नहीं; \omterm{def:g2:animals-grow-up:born}{अंडा} दूसरा \omterm{def:g6:cell-unit-of-life:parts}{केंद्रक}
\emph{और} भंडार-कोठरी, दोनों लाया था। और आरोपण, बाक़ी बातों के अलावा,
भंडार-कोठरी के चुकते जाने और किरायेदार के कमरे की सेवा से जुड़ जाने का
नाम है।
\end{example}

\section{भ्रूण, नाल, गर्भस्थ शिशु}

\begin{definition}[नाल]\label{def:g8:from-egg-to-birth:placenta}
आरोपित गेंद से एक ही साथ दो रचनाएँ उगती हैं: ख़ुद \emph{\omterm{def:g5:human-reproduction:pregnancy}{भ्रूण}}, और
उसके जीवन-आधार का जोड़ --- यानी \emph{नाल}\index{नाल}, \omterm{prop:g5:human-reproduction:plan}{गर्भाशय} की
दीवार में उगी एक तश्तरी, जहाँ माँ का \omterm{prop:g4:journey-of-food:absorb}{रक्त} और \omterm{def:g5:human-reproduction:pregnancy}{भ्रूण} का \omterm{prop:g4:journey-of-food:absorb}{रक्त} \emph{बिना
मिले} बहुत क़रीब आ जाते हैं। उसकी विशाल पतली सतहों के पार ---
\cref{prop:g7:how-animals-breathe:spec} वाली विशेष-सूची, एक बार फिर बनी
हुई --- \omterm{def:g7:digestion-and-nutrients:families}{पोषक तत्व} और \omterm{def:g7:what-is-respiration:respiration}{ऑक्सीजन} \omterm{def:g5:human-reproduction:pregnancy}{भ्रूण} की ओर जाते हैं, और कचरा तथा \omterm{def:g7:what-is-respiration:respiration}{कार्बन
डाइऑक्साइड} वापस आते हैं (\cref{prop:g5:human-reproduction:cord} वाली
रज्जु नाल और \omterm{def:g5:human-reproduction:pregnancy}{भ्रूण} के बीच चलती है)। उस मुसाफ़िर के लिए नाल ही \omterm{def:g7:how-animals-breathe:four}{फेफड़ा} है,
आँत है और \omterm{def:g7:removing-wastes:kidneys}{गुर्दा} भी --- यानी तीनों सरहदें एक ही अंग में।
\end{definition}

\begin{proposition}[समय-सारणी, साज़ों के साथ दोबारा]\label{prop:g8:from-egg-to-birth:timetable}
\Cref{ex:g5:human-reproduction:timetable} वाले पड़ाव, अब समझाए हुए:
\begin{itemize}
  \item हफ़्ते 1--8, यानी \emph{\omterm{def:g5:human-reproduction:pregnancy}{भ्रूण}}: शरीर का नक़्शा बिछाया जाता है
        --- \omterm{def:g5:first-look-at-cells:cell}{कोशिकाएँ} बँटती हैं, जगह बदलती हैं और
        \emph{विशेषज्ञ बनती} हैं
        (\cref{rem:g6:cell-unit-of-life:twoways}) --- \omterm{def:g4:heart-and-blood:heart}{दिल},
        \omterm{def:g5:brain-in-command:system}{तंत्रिका}, \omterm{def:g3:bones-and-muscles:skeleton}{हड्डी} और बाक़ी सब में; और \omterm{def:g4:heart-and-blood:heart}{दिल} चौथे हफ़्ते से ही
        धड़कता है, क्योंकि भंडार-और-नाल वाली इस अर्थव्यवस्था को अपना
        पंप जल्दी चाहिए;
  \item हफ़्ते 9--40, यानी \emph{\omterm{def:g5:human-reproduction:pregnancy}{गर्भस्थ शिशु}}: नक़्शा बन चुका, अब
        बढ़त और पकाई --- \omterm{def:g5:brain-in-command:system}{मस्तिष्क} ख़ुद को जोड़ता हुआ, \omterm{def:g7:how-animals-breathe:four}{फेफड़े} ख़ाली ही
        रिहर्सल करते हुए
        (\cref{exo:g5:human-reproduction:11} ने उनके पहले रोने वाले
        पदार्पण की बात कही थी), और चौथे महीने की लातें काम करती
        \omterm{def:g3:bones-and-muscles:muscle}{पेशियों} का ऐलान करती हुईं;
  \item और पूरे दौरान: \omterm{def:g8:puberty:hormones}{हॉर्मोन} \omterm{prop:g5:human-reproduction:plan}{गर्भाशय} को शांत और चक्र को ठहरा हुआ
        थामे रहते हैं --- यानी \cref{def:g8:puberty:hormones} वाला
        रसायन, \omterm{def:g5:human-reproduction:pregnancy}{गर्भावस्था} का दफ़्तर चलाता हुआ।
\end{itemize}
\end{proposition}
\admitted

\begin{example}[शुरुआती हफ़्ते सबसे ज़्यादा देखभाल क्यों माँगते हैं]\label{ex:g8:from-egg-to-birth:earlycare}
\Cref{prop:g5:human-reproduction:care} ने चेताया था कि तंबाकू और शराब
शिशु तक पार कर जाते हैं; और समय-सारणी बताती है कि यह सबसे ज़्यादा कब
मायने रखता है: नक़्शा बिछाने वाले \omterm{def:g5:human-reproduction:pregnancy}{भ्रूण} के हफ़्तों में, जब हर अंग
विशेषज्ञ बनती \omterm{def:g5:first-look-at-cells:cell}{कोशिकाओं} की एक छोटी-सी भीड़ होता है, तब एक भी बिगड़ा हुआ
निर्देश पूरी बनावट में गूँज सकता है --- और ठीक वही हफ़्ते हैं जिनमें
\omterm{def:g5:human-reproduction:pregnancy}{गर्भावस्था} का पता तक नहीं चलता। इसीलिए डॉक्टरों का वह सीधा-सादा गणित:
देखभाल पहली जाँच से नहीं, पहली संभावना से शुरू होती है।
\end{example}

\section{जन्म, और सरहदों की अदला-बदली}

\begin{proposition}[जन्म, हॉर्मोन के हुक्म से हुई अदला-बदली]\label{prop:g8:from-egg-to-birth:birth}
चालीसवें हफ़्ते के आस-पास हॉर्मोनों के इशारे --- शिशु के और माँ के,
आपस में बात करते हुए --- \omterm{prop:g5:human-reproduction:plan}{गर्भाशय} की उस लंबी शांति को ख़त्म कर देते हैं:
\begin{enumerate}
  \item \emph{प्रसव-पीड़ा}: \omterm{prop:g5:human-reproduction:plan}{गर्भाशय}, यानी शरीर की सबसे ताक़तवर \omterm{def:g3:bones-and-muscles:muscle}{पेशी}
        (\cref{rem:g8:reproductive-systems:living} उसके नरम काम से
        मिल चुकी है), बढ़ती हुई लहरों में सिकुड़ती है, रास्ता खोलती है
        और शिशु को जन्म की नली में नीचे धकेलती है;
  \item \emph{जन्म}: आम तौर पर \omterm{def:g1:my-body:parts}{सिर} पहले
        (\cref{ex:g5:human-reproduction:birth});
  \item \emph{अदला-बदली}: पहली रुलाई कभी इस्तेमाल न हुए \omterm{def:g7:how-animals-breathe:four}{फेफड़े} भर देती
        है; रज्जु की आवाजाही ख़त्म होती है और उसे काट दिया जाता है; और
        मिनटों में नवजात अपने ही फेफड़ों, आगे चलने वाली आँत और \omterm{def:g7:removing-wastes:kidneys}{गुर्दों}
        पर चलने लगता है --- यानी वे सारी सरहदें, जो अब तक \omterm{def:g8:from-egg-to-birth:placenta}{नाल} सँभाल
        रही थी (\cref{exo:g5:human-reproduction:11}, अब पूरे नक़्शे के
        साथ);
  \item \emph{आँवल}: \omterm{def:g8:from-egg-to-birth:placenta}{नाल}, अपना काम पूरा करके, अलग होती है और
        पीछे-पीछे निकल आती है --- यानी वह अकेला अंग, जिसे शरीर एक ही
        इस्तेमाल के लिए बनाता है और फिर छोड़ देता है।
\end{enumerate}
\end{proposition}
\admitted

\begin{example}[नवजात की पहली जाँच-पड़ताल]\label{ex:g8:from-egg-to-birth:audit}
दाइयाँ जन्म के मिनटों बाद ही नवजात को अंक देती हैं --- साँस, \omterm{def:g4:heart-and-blood:heart}{दिल} की
धड़कन, रंग, हरकत, रुलाई। इस जाँच-सूची को इस किताब की \omterm{def:g1:the-five-senses:sense}{आँखों} से पढ़ो: यह
तो \cref{prop:g7:removing-wastes:complete} वाला आपूर्ति-और-हटाव का
फंदा है, जिसकी उसकी पहली अकेली दौड़ पर जाँच हो रही है। हर इंसान सबसे
पहले जो इम्तिहान पास करता है वह शरीर-क्रिया का इम्तिहान है।
\end{example}

\begin{method}[गर्भावस्था की किसी भी अवस्था को समझाना]\label{met:g8:from-egg-to-birth:explain}
उन नौ महीनों के किसी भी पल के बारे में पूछा जाए, तो चार साज़ों से जवाब
दो:
\begin{enumerate}
  \item \emph{\omterm{def:g5:first-look-at-cells:cell}{कोशिकाएँ}}: उनका बँटना और विशेषज्ञ बनना क्या कर रहा है?
  \item \emph{आपूर्ति}: भंडार-कोठरी, या \omterm{def:g8:from-egg-to-birth:placenta}{नाल} --- क्या पार होता है, और
        क्या हरगिज़ नहीं?
  \item \emph{\omterm{def:g8:puberty:hormones}{हॉर्मोन}}: क्या थामे रखा जा रहा है, या क्या हुक्म हो रहा
        है?
  \item \emph{देखभाल}: यह अवस्था माँ की आदतों और डॉक्टरों के कैलेंडर
        से क्या माँगती है?
\end{enumerate}
\end{method}

\section{अभ्यास}

\begin{exercise}[$\star$]\label{exo:g8:from-egg-to-birth:1}
पहले हफ़्ते का समय बताओ: बँटने की शुरुआत, शहतूत जैसी गेंद, और आरोपण।
\end{exercise}

\begin{exercise}[$\star$]\label{exo:g8:from-egg-to-birth:2}
पहले हफ़्ते का मुसाफ़िर किस पर जीता है, और वह इंतज़ाम कौन ख़त्म करता
है?
\end{exercise}

\begin{exercise}[$\star$]\label{exo:g8:from-egg-to-birth:3}
\omterm{def:g8:from-egg-to-birth:placenta}{नाल} क्या है --- उसके पार हर तरफ़ क्या-क्या जाता है, और क्या कभी नहीं
जाता?
\end{exercise}

\begin{exercise}[$\star$]\label{exo:g8:from-egg-to-birth:4}
चालीस हफ़्तों को \omterm{def:g5:human-reproduction:pregnancy}{भ्रूण} और \omterm{def:g5:human-reproduction:pregnancy}{गर्भस्थ शिशु} में बाँटो, और हर चरण का काम भी
बताओ।
\end{exercise}

\begin{exercise}[$\star$]\label{exo:g8:from-egg-to-birth:5}
सारे अंगों में \omterm{def:g4:heart-and-blood:heart}{दिल} ही चौथे हफ़्ते से पहले-पहल क्यों धड़कता है?
\end{exercise}

\begin{exercise}[$\star$]\label{exo:g8:from-egg-to-birth:6}
जन्म की चारों गिनी हुई घटनाएँ क्रम से गिनाओ।
\end{exercise}

\begin{exercise}[$\star\star$]\label{exo:g8:from-egg-to-birth:7}
\omterm{def:g8:from-egg-to-birth:placenta}{नाल} किन तीन अंगों का काम करती है? हर एक की सरहद वाले अध्याय का नाम भी
बताओ।
\end{exercise}

\begin{exercise}[$\star\star$]\label{exo:g8:from-egg-to-birth:8}
\omterm{def:g5:human-reproduction:pregnancy}{भ्रूण} वाले हफ़्ते सबसे ज़्यादा देखभाल क्यों माँगते हैं ---
\cref{ex:g8:from-egg-to-birth:earlycare} में से दो कारण?
\end{exercise}

\begin{exercise}[$\star\star$]\label{exo:g8:from-egg-to-birth:9}
नवजात की अंक-सूची को \cref{prop:g7:removing-wastes:complete} वाले
फंदे की तरह पढ़ो: साँस, \omterm{def:g4:heart-and-blood:heart}{दिल} की धड़कन और रंग को उनके तंत्रों से मिलाओ।
\end{exercise}

\begin{exercise}[$\star\star$]\label{exo:g8:from-egg-to-birth:10}
जुड़वाँ, फिर से (\cref{rem:g5:human-reproduction:twins}):
\cref{prop:g8:from-egg-to-birth:firstweek} के किस पल पर वह एक \omterm{prop:g8:sexual-reproduction:core}{युग्मनज}
दो में बँटता है, और उस समय की वजह से वे ``नक़लों जैसे एक-से'' क्यों
होते हैं?
\end{exercise}

\begin{exercise}[$\star\star$]\label{exo:g8:from-egg-to-birth:11}
\cref{met:g8:from-egg-to-birth:explain} को चौथे महीने पर चलाओ ---
यानी पहली लातों वाले महीने पर।
\end{exercise}

\begin{exercise}[$\star\star\star$]\label{exo:g8:from-egg-to-birth:12}
``जन्म सरहदों का बदल जाना है।'' इसे सही ठहराओ: वे सारी अदला-बदलियाँ
गिनाओ जो \omterm{def:g8:from-egg-to-birth:placenta}{नाल} सँभाल रही थी, और वह अंग भी जो पहली रुलाई पर हर एक अपने
ज़िम्मे लेता है --- और साथ में वह मिनटों का समय भी, जो इस अदला-बदली को
उल्लेखनीय बना देता है।
\end{exercise}

\section{समस्या: दाई की समय-रेखा}

\begin{problem}[{सप्ताहांत समस्या --- एक गर्भावस्था, मुलाक़ात-दर-मुलाक़ात
देखी हुई}]\label{pb:g8:from-egg-to-birth:1}
एक (कल्पित) \omterm{def:g5:human-reproduction:pregnancy}{गर्भावस्था} की दाई वाली फ़ाइल: आठवें हफ़्ते में पहली दर्ज
मुलाक़ात, बारहवें और बीसवें हफ़्ते में अल्ट्रासाउंड, अट्ठाईसवें तथा
छत्तीसवें में जाँच, और चालीसवें हफ़्ते में जन्म। इस फ़ाइल को अध्याय के
साज़ों से पढ़ो।

\medskip\noindent\textbf{भाग I --- फ़ाइल के शुरुआती पन्ने।}
\begin{enumerate}
  \item \omterm{def:g5:human-reproduction:pregnancy}{गर्भावस्था} की पुष्टि पहली मुलाक़ात से हफ़्तों पहले एक
        हॉर्मोन-जाँच से हुई थी। ऐसी जाँच कौन-सा इशारा पढ़ती है, और
        \cref{prop:g8:from-egg-to-birth:firstweek} की किस घटना के बाद
        से वह मौजूद होता है?
  \item आठवें हफ़्ते पर दाई लिखती है ``\omterm{def:g5:human-reproduction:pregnancy}{भ्रूण} नक़्शे के हिसाब से पूरा,
        लगभग सेम के दाने जितना''। इसे समय-सारणी की भाषा में उतारो ---
        यहाँ कौन-सा चरण बंद होता है?
  \item पहली दर्ज मुलाक़ात की सलाह सबसे ऊपर यही रखती है: कोई शराब
        नहीं, कोई धुआँ नहीं --- और फ़ोलिक अम्ल, यानी उसारी वाला एक
        विटामिन, ``बेहतर हो तो शुरुआत से भी पहले से''। तीनों का औचित्य
        \cref{ex:g8:from-egg-to-birth:earlycare} से दो।
  \item बारहवें हफ़्ते का अल्ट्रासाउंड वह धड़कन दिखाता है जो चौथे
        हफ़्ते में शुरू हुई थी। आपूर्ति की भाषा में, इतनी जल्दी
        क्यों?
\end{enumerate}

\medskip\noindent\textbf{भाग II --- बीच के महीने।}
\begin{enumerate}[resume]
  \item बीसवें हफ़्ते का अल्ट्रासाउंड हर अंग और \omterm{def:g8:from-egg-to-birth:placenta}{नाल} की जगह नापता है।
        \omterm{def:g5:human-reproduction:pregnancy}{गर्भस्थ शिशु} वाले चरण का काम बताओ, और \omterm{def:g8:from-egg-to-birth:placenta}{नाल} के वे तीनों काम भी
        जिनकी जाँच हो रही है।
  \item माँ अठारहवें हफ़्ते के आस-पास पहली लातों की ख़बर देती है।
        कौन-सी मशीनरी अपना ऐलान कर रही है
        (\cref{prop:g5:brain-in-command:loop} उस कड़ी का नाम बताती
        है)?
  \item अट्ठाईसवें हफ़्ते पर दाई माँ के \omterm{prop:g4:journey-of-food:absorb}{रक्त} की शर्करा और दबाव जाँचती
        है। माँ का रक्त-चक्र दोहरे बोझ में क्यों है --- वह किन-किन की
        सरहदें चला रहा है?
  \item हफ़्ते 30--36: \omterm{def:g5:human-reproduction:pregnancy}{गर्भस्थ शिशु} मुख्य रूप से बढ़ता है और \omterm{def:g7:how-animals-breathe:four}{फेफड़े}
        रिहर्सल करते हैं। रिहर्सल किस चीज़ की --- और वह भी किस अकेले
        पल के लिए?
\end{enumerate}

\medskip\noindent\textbf{भाग III --- आख़िरी पन्ना।}
\begin{enumerate}[resume]
  \item चालीसवाँ हफ़्ता: प्रसव-पीड़ा। उस \omterm{def:g3:bones-and-muscles:muscle}{पेशी} का नाम बताओ, उसके दोनों
        काम, और यह भी कि हुक्म किसने दिया
        (\cref{prop:g8:from-egg-to-birth:birth})।
  \item फ़ाइल की जन्म वाली प्रविष्टि: ``फ़ौरन ज़ोरदार रुलाई; रज्जु
        बाँधकर काटी गई; \omterm{def:g8:from-egg-to-birth:placenta}{नाल} पूरी निकली।'' इन तीनों उपवाक्यों की जाँच
        अदला-बदली से करो।
  \item नवजात मिनटों वाली जाँच-सूची पर ऊँचे अंक पाता है।
        \cref{ex:g8:from-egg-to-birth:audit} की भाषा में बताओ कि वह
        जाँच-सूची किस बात की गवाही देती है।
  \item फ़ाइल को माता-पिता के लिए एक वाक्य से बंद करो: उन चालीस
        हफ़्तों ने किस अकेली \omterm{def:g5:first-look-at-cells:cell}{कोशिका} से, किस मशीनरी के ज़रिए, क्या बना
        दिया।
\end{enumerate}
\end{problem}
